\newcommand{\uniConverge}[2]{\toInf  \sup_{#1} \left| {#2} \right| = 0}
\newcommand{\estLike}{Q \left( \est_{n} \mid \paramVal \right) }

Many probabilistic models are naturally formulated in terms of a sequence of random variables whose length is not
fixed in advance.
For example, when a coin is tossed repeatedly, there is no natural upper bound on the number of tosses, and any
finite sample can be viewed as an initial segment of an infinite sequence.

Motivated by such examples, let $S_n$ denote the first $n$ variables of an
infinite sequence of random variables:
\[
    S_n = (X_1, X_2, \ldots, X_n).
\]
The collection \infSeq{S} therefore forms a sequence of random variables.

In many settings, our uncertainty about $S_n$ can be represented by a central distribution.
A prominent example arises when the sequence is assumed to be exchangeable.
As discussed earlier, exchangeability can be viewed as a form of partial knowledge that defines a set of plausible
distributions, from which a central distribution for $S_n$ can be constructed.

However, this central distribution generally depends on the sequence length $n$.
According to the principle of objectivity, our beliefs should depend only on the available information and not on
arbitrary aspects of the experimental setup, such as the number of variables under consideration.

Importantly, a distribution defined over a longer sequence induces marginal distributions over all shorter
initial segments.
This suggests that, rather than constructing a separate central distribution for each
$S_n$, we should seek a single probability measure over the infinite sequence, from which the distributions of all
finite initial segments can be obtained as marginals.

This motivates the construction of a central distribution that is well defined over $S_{\infty}$.

\begin{definition}
    \label{def:infseq-central}
    Let \infSeq{S} be a sequence of random variables and \setQ be a set of probability measures,
    where each element of \setQ defines a distribution over each $S_n$.
    Assume \cDist[n] is a central distribution for $S_n$.
    Let \cDist be a probability measure that defines a distribution over every $S_n$.
    \cDist is a central distribution of \setQ for \infSeq{S}, if we have
    \begin{equation*}
        \toInf \kl[{\cDist[n]}][\cDist][S_n] = 0.
    \end{equation*}
\end{definition}

Intuitively, this definition states that as $n \to \infty$ the distribution that $C$ induces over $S_n$ becomes
increasingly similar to a central distribution for $S_n$.
In other words, \cDist is a central distribution for the infinite sequence $S_{\infty}$.

Our next goal is to extend Theorem~\ref{th:lambda-centrality} so that it also applies in this asymptotic setting,
thereby providing a characterization of central distributions for infinite sequences.

\begin{theorem}[Asymptotic Centrality]
    \label{th:limiting-centrality}
    Let \infSeq{S} be a sequence of random variables and \setQ be a set of probability measures,
    where each element of \setQ defines a distribution over each $S_n$.
    Assume that there exists a natural number $M$ such that,
    for every $n \geq M$, there exists a distribution $C_n$ satisfying the
    conditions of Theorem~\ref{th:lambda-centrality} for $S_n$.

    Let \cDist be a probability measure that defines a distribution over every $S_n$.
    If there exists a sequence of constants \infSeq{\lambda} such that
    \[
        \uniConverge{\inSet{Q}}{\kl[][\cDist][S_n] - \lambda_n},
    \]
    then \cDist is a central distribution of \setQ for \infSeq{S}.
\end{theorem}

\begin{proof}
    For every $n \geq M$, by Theorem~\ref{th:lambda-centrality}, there exists a constant \(\mu_n\) such that for
    all \inSet{Q}
    \[
        \kl[][\cDist_n][S_n] = \mu_n .
    \]
    By the definition of a central distribution we have
    \begin{equation*}
        \sup_{\inSet{Q}} \kl[][\cDist][S_n] \geq \mu_n,
        \label{eq:th_lim_central_proof_sup}
    \end{equation*}
    and, by Lemma~\ref{lm:imperfection}, it also follows that
    \begin{equation*}
        \inf_{\inSet{Q}} \kl[][\cDist][S_n] \leq \mu_n .
        \label{eq:th_lim_central_proof_inf}
    \end{equation*}

    On the other hand, clearly
    \[
        \left| \sup_{\inSet{Q}} \kl[][\cDist][S_n] - \lambda_n \right|
        \leq \sup_{\inSet{Q}}  \left|\kl[][\cDist][S_n] - \lambda_n \right|,
    \]
    and,
    \[
        \left|\inf_{\inSet{Q}} \kl[][\cDist][S_n] - \lambda_n \right|
        \leq \sup_{\inSet{Q}}  \left|\kl[][\cDist][S_n] - \lambda_n \right|.
    \]

    Consequently, by the squeeze theorem, the sequence \infSeq[]{\mu_n - \lambda_n} must converge to $0$,
    and for any arbitrary \(Q \in \setQ\), we have
    \begin{align*}
        \toInf \sum_{s_n} Q(s_n) \ln \frac{\cDist_n(s_n)}{\cDist(s_n)} &=
        \toInf (\kl[][\cDist][S_n] - \kl[][\cDist_n][S_n]) \\
        &= \toInf (\kl[][\cDist][S_n] -\mu_n + \lambda_n - \lambda_n) \\
        &= \toInf(\kl[][\cDist][S_n] - \lambda_n) - \toInf(\mu_n - \lambda_n) \\
        &= 0 .
    \end{align*}
    Since $\cDist_n$ satisfies the conditions of Theorem~\ref{th:lambda-centrality}, it is a convex combination of
    elements of \setQ.
    Therefore, $\kl[\cDist_n][\cDist][S_n]$ can be expressed as a convex combination of the quantities
    \[
        Q_i(s_n) \ln \frac{\cDist_n(s_n)}{\cDist(s_n)},
    \]
    where $Q_i \in \setQ$.

    By the previous result, we can conclude
    \[
        \toInf \kl[\cDist_n][\cDist][S_n] = 0.
    \]
    Thus, by Definition~\ref{def:infseq-central}, \cDist is a central distribution of \setQ for the
    sequence \infSeq{S}.
\end{proof}

We now present an asymptotic version of Theorem~\ref{th:estimator-to-var}, applicable to infinite sequences of
random variables.

\begin{theorem}
    \label{th:estimator-to-infseq}
    Let \infSeq{S} and \infSeq{T} be two sequences of random variables.
    Let \setQ be a set of probability measures, where each \inSet{Q} defines a joint distribution
    over $(S_n, T_n)$ for every $n$.

    Assume there exists a natural number $M$ and a fixed conditional distribution $R$
    such that, for all \inSet{Q} and all $n \geq M$,
    \[
        Q(S_n, T_n) = R(S_n \mid T_n)Q(T_n).
    \]

    Let $C$ be a central distribution of \setQ for \infSeq{T}.
    Define
    \[
        C^*(S_n, T_n) = R(S_n \mid T_n)C(T_n), \quad n \geq M.
    \]
    If $C^*$ is a probability measure, then it is a central distribution of \setQ for \infSeq[]{(S_n, T_n)}.
\end{theorem}

\begin{proof}
    By Theorem~\ref{th:estimator-to-var}, for each $n \geq M$ we have
    \[
        C_n(S_n, T_n) = R(S_n \mid T_n) C_n(T_n),
    \]
    where $C_n(T_n)$ is a central distribution for $T_n$, and $C_n(S_n,  T_n)$ is a central distribution for
    $(S_n, T_n)$.
    Given that $C^*$ is a probability measure, we have
    \begin{align*}
        \kl[{C_n}][C^*][S_n, T_n] &= \kl[{R(S_n \mid T_n)C_n}][R(S_n \mid T_n)C][T_n] \\
        &= \kl[{C_n}][C][T_n].
    \end{align*}
    Since $C$ is a central distribution for \infSeq{T}, by definition,
    \[
        \toInf \kl[{C_n}][C][T_n] = 0.
    \]
    It follows that
    \[
        \toInf \kl[{C_n}][C^*][S_n, T_n] = 0,
    \]
    and therefore $C^*$ is a central distribution for \infSeq[]{(S_n, T_n)}.
\end{proof}

The following construction is heuristic and is intended to motivate the asymptotic form of central priors under
standard regularity conditions.
A complete rigorous treatment is left for future work.

Let \setQ be a family of probability distributions indexed by a parameter vector \param.
Suppose \infSeq{\est} is a sequence of consistent estimators of \param.
Let $P(\paramVal \mid\est_n)$ denote the posterior distribution obtained from an arbitrary smooth positive prior.

Under regular parametric conditions, posterior distributions based on consistent estimators become asymptotically
insensitive to the choice of smooth positive prior.

Motivated by this phenomenon, we heuristically derive the asymptotic form of a central prior by considering the
posterior density evaluated at the point $t=\paramVal$.
A similar asymptotic construction also appears in~\cite{BERNARDO200517}.

Specifically, we define
\begin{equation}
    \label{eq:central-limit}
    \cPrior(\paramVal) \propto \toInf \frac{P(\paramVal \mid\est_n = t)|_{t = \paramVal}}{\mu_n},
\end{equation}
where \infSeq{\mu} is a sequence of positive constants chosen so that the limit exists and the convergence is
uniform over \paramVal.

Let $k$ denote the normalizing constant converting proportionality into equality.
Since the logarithm is a continuous function, Equation~\ref{eq:central-limit} implies
\begin{equation*}
    \uniConverge{\paramVal}{\ln \frac{P(\paramVal \mid\est_n = t)|_{t = \paramVal}}
    {\cPrior(\paramVal)} - \ln \mu_n + \ln k }.
\end{equation*}

For every \paramVal, since $\est_n$ is a consistent estimator, the sampling distribution $\estLike$ becomes
increasingly concentrated around $\paramVal$ as $n \to \infty$.
Therefore, we may write
\begin{equation*}
    \uniConverge{\paramVal}{\sum_{t_n} Q(t_n \mid \paramVal) \ln \frac{P(\paramVal \mid t_n)}
    {\cPrior(\paramVal)} - \ln \mu_n + \ln k }.
\end{equation*}

Furthermore, under the usual regularity conditions, the posterior distribution becomes asymptotically
insensitive to the choice of smooth positive prior.
Motivated by this asymptotic prior-insensitivity, we heuristically replace $P(\paramVal \mid \est_n)$
with $c(\paramVal \mid \est_n)$,
obtaining
\begin{equation*}
    \uniConverge{\paramVal}{ \sum_{t_n} Q(t_n \mid \paramVal) \ln \frac{Q(t_n \mid \paramVal)}{C(t_n)} - \ln \mu_n + \ln k }.
\end{equation*}
Letting
\[
    \lambda_n = \ln \mu_n - \ln k,
\]
we observe that
\begin{equation*}
    \uniConverge{\paramVal}{\Dkl{\estLike}{C(\est_n)} - \lambda_n}.
\end{equation*}
Under appropriate additional assumptions, the conditions of Theorem~\ref{th:limiting-centrality} are satisfied.
Thus, under the stated regularity assumptions and heuristic asymptotic arguments, $\cPrior(\paramVal)$ is
expected to act as a central prior for the estimator sequence \infSeq{\est}.

A central distribution over an estimator sequence cannot by itself be used to perform inference at the level of observed
outcomes.
To do so, we must convert it into a distribution over the primary variables of interest.
Theorems~\ref{th:estimator-to-var} and~\ref{th:estimator-to-infseq} provide a simple
method for this conversion.

When dealing with layered ignorance and seeking the central distribution of a hierarchy of distributions, as defined
in Definition~\ref{def:hierarchy-of-distributions}, we typically encounter situations in which a given level of the
hierarchy contains only finitely many alternatives.
In such cases, and under appropriate regularity conditions, if there exists a consistent estimator for the index
identifying these alternatives, then the central distribution reduces to an equally weighted average of the
corresponding representative central distributions.

\begin{example}[Binomial Experiment]
    \newcommand{\pv}{\paramVal}
    Consider a binomial experiment consisting of $n$ independent Bernoulli trials.
    Let \setQ denote the family of probability distributions corresponding to different values of the success
    probability $\pv \in [0,1]$, with \pv serving as the indexing parameter for \setQ.
    Let $T_n$ be the sample proportion of successes in the $n$ trials.
    Since $T_n$ is a consistent estimator of \pv, we can apply Equation~\ref{eq:central-limit} to determine a
    central distribution for the
    sequence \infSeq{T}.

    Assuming a uniform prior on $\pv$, Bayes’ rule gives the posterior distribution
    \[
        P\left( \pv \, \middle| \, T_n = \tfrac{k}{n} \right) = (n+1)\binom{n}{k} \pv^k (1 - \pv)^{n - k},
    \]
    which is a $\mathrm{Beta}(k+1, n-k+1)$ distribution.
    Using Stirling’s formula, we derive its asymptotic form:
    \begin{equation*}
        P\left( \pv \, \middle| \, T_n = \tfrac{k}{n} \right) \sim
        \frac{(n+1)n^{(n+\hf)}}{\sqrt{2 \pi} (n-k)^{(n-k+\hf)} k^{(k+\hf)}} \, \pv^k (1 - \pv)^{n - k}.
    \end{equation*}
    Evaluating the asymptotic expression at $\frac{k}{n} = \pv$, we obtain
    \begin{equation*}
        P(\paramVal \mid\est_n = t)|_{t = \paramVal} \sim \frac{n+1}{\sqrt{2 \pi n}} \, \pv^{-\hf}(1 - \pv)^{-\hf}.
    \end{equation*}
    Therefore, if we define
    \[
        \mu_n = \frac{n+1}{\sqrt{2 \pi n}},
    \]
    the following limit exists and the convergence is uniform over $\paramVal \in (0,1)$:
    \[
        \toInf \frac{P(\paramVal \mid\est_n = t)|_{t = \paramVal}}{\mu_n} = \pv^{-\hf}(1 - \pv)^{-\hf}.
    \]
\end{example}

\begin{example}[Fused Multinomial Model]
    \label{ex:fused-multinomial}
    \newcommand{\prodm}{\prod_{y=1}^m}
    \newcommand{\Qt}{Q_{\theta}}
    Let \infSeq{S} be a sequence in which each $S_n = (\seqExp{X})$ represents $n$ independent trials of an
    experiment with $k$ possible outcomes,
    and \setQ be a set of probability distributions over $S_n$ indexed by a real vector \paramVal.
    Let $f: \set{X} \to \set{Y}$ be a deterministic function that maps each outcome of $X_i$
    to a label $Y_i = f(X_i)$, where $\set{Y} = \{1, \dots, m\}$.
    Assume that for all $i$ and $\Qt \in \setQ$, if $f(x_i) = f(x'_i)$ we have \[\Qt(\seqExp{x}) = \Qt(\seqExp{x'}).\]
    That is, the conditional distribution $\Qt(X_i \mid Y_i = y)$ is uniform over the preimage $f^{-1}(y)$.

    We can consider $\theta = (\theta_1, \dots, \theta_m)$ to be a probability vector over $\mathcal{Y}$,
    where each $\theta_y$ specifies the probability of observing label $y$.
    For a sequence $S_n$, let $n_y$ denote the number of times label $y \in \mathcal{Y}$ appears.
    Let $|f^{-1}(y)|$ denote the number of outcomes in $\mathcal{X}$ that are mapped to label $y$ by $f$.
    Then the probability of observing $S_n$ under $\Qt \in \setQ$, denoted $Q(S_n \mid \theta)$, is:
    \[
        Q(S_n \mid \theta) = \prodm \left( \frac{1}{|f^{-1}(y)|} \cdot \theta_y \right)^{n_y}.
    \]

    Now, consider the estimator sequence $T_n = \frac{1}{n} \left( \seqExp[m]{n} \right)$, which is a consistent
    estimator for $\theta$.
    The conditional distribution $\Qt(S_n \mid T_n)$ is then given by:
    \begin{equation}
        R(S_n \mid T_n) = \frac{1}{n!} \prodm \frac{n_y!}{|f^{-1}(y)|^{n_y}}\,.
        \label{eq:fused-multinomial-conditional}
    \end{equation}
    Note that this distribution is independent of $\theta$.
    Therefore, by Theorem~\ref{th:estimator-to-infseq}, a central distribution for $S_n$ can be obtained from a
    central distribution for $T_n$.

    As $n \to \infty$, the distribution of $T_n$ converges to a multivariate Gaussian distribution.
    Using this fact and Equation~\ref{eq:central-limit}, we can determine a central prior for \infSeq{T}:
    \[
        c(\theta) \propto \prodm \theta_y^{-\hf}.
    \]
    This is a Dirichlet prior and induces a Dirichlet-multinomial distribution for each $T_n$:
    \[
        C(T_n) = \frac{\gm{\hf[m]}\gm{n+1}}{\gm{n + \hf[m]}}
        \prodm \frac{\gm{n_y + \hf}}{\gm{\hf}\gm{n_y+1}}\,.
    \]
    Note that $C$ is a central distribution for the sequence \infSeq{T}, viewed as a probability measure on the
    infinite sequence $T_\infty$.
    The notation $C(T_n)$ denotes the marginal distribution of $T_n$.

    By Theorem~\ref{th:estimator-to-infseq} and Equation~\ref{eq:fused-multinomial-conditional}, we can derive a
    central distribution for \infSeq{S}, which induces the following distribution over each $S_n$:
    \[
        C(S_n) = \frac{\gm{\hf[m]}}{\gm{n + \hf[m]}}
        \prodm \frac{\gm{n_y + \hf}}{\gm{\hf} |f^{-1}(y)|^{n_y}}\,.
    \]
\end{example}
